\chapter{Невозможность получить отображение момента из обычного спектрального функционала}

Внешний разбор указал на повторяющуюся попытку вывести ориентированную
динамику колчана из неориентированного спектрального следа. Здесь доказывается
точное ограниченное утверждение для функционалов, зависящих только от спектра
обычного нечётного оператора Дирака. Оно не относится ко всей
некоммутативной геометрии.

\section{Ориентированная трёхвершинная цепочка}

Пусть
\begin{equation}
 \mathcal H=\mathcal H_0^+\oplus\mathcal H_1^-\oplus\mathcal H_2^+,
 \qquad
 D=\begin{pmatrix}0&X^\dagger&0\\X&0&Y^\dagger\\0&Y&0\end{pmatrix}.
 \label{eq:v5-three-node-dirac}
\end{equation}
С оператором
\begin{equation}
 M=(X\;\;Y^\dagger):\mathcal H_0\oplus\mathcal H_2\to\mathcal H_1
\end{equation}
имеем \(D=\left(\begin{smallmatrix}0&M^\dagger\\M&0\end{smallmatrix}\right)\).
Обозначим
\begin{equation}
 A=XX^\dagger,\qquad B=Y^\dagger Y.
\end{equation}
Тогда \(MM^\dagger=A+B\), тогда как ориентированное отображение момента
равно \(\mu=A-B\).

\section{Теорема о спектральной неразличимости}

Ненулевые спектры \(M^\dagger M\) и \(MM^\dagger\) совпадают. Поэтому для
любой функции \(f\), определённой на конечном спектре,
\begin{equation}
 \Tr f(D^2)=2\Tr_{\mathcal H_1}f(A+B)+c_f,
 \label{eq:v5-general-spectral-functional}
\end{equation}
где \(c_f\) зависит только от числа нулевых мод и \(f(0)\).

\begin{theorem}[Спектральная слепота к входящим и исходящим стрелкам]
Любой обычный унитарно-инвариантный функционал, построенный только из
спектра \(D^2\), зависит от стрелок \(X,Y\) через сумму \(A+B\) и не может
восстановить \(\Tr(A-B)^2\) на всём пространстве конфигураций.
\end{theorem}

\begin{proof}
Зафиксируем положительный оператор \(S\) и \(0\le t\le1\). Положим
\(A_t=tS\), \(B_t=(1-t)S\). Такие операторы допускают требуемые разложения
через \(X_t,Y_t\). Сумма \(A_t+B_t=S\) не зависит от \(t\), поэтому все
спектральные функционалы постоянны, но
\begin{equation}
 \Tr(A_t-B_t)^2=(2t-1)^2\Tr S^2
\end{equation}
меняется вместе с \(t\).
\end{proof}

\section{Квартичный частный случай}

\begin{align}
 \Tr D^4&=2\Tr(A+B)^2
 =2\Tr A^2+2\Tr B^2+4\Tr AB,
 \label{eq:v5-positive-cross}\\
 \Tr\mu^2&=\Tr A^2+\Tr B^2-2\Tr AB.
 \label{eq:v5-negative-cross}
\end{align}
При положительном коэффициенте перед \(\Tr D^4\) смешанный член имеет знак,
противоположный отображению момента. Этим объясняются обе уже найденные
картины: \((\rho^2+r^2)^2\) вместо \((\rho^2-r^2)^2\) и исчезновение
триплетных связующих полей в активном прямоугольнике.

\section{Почему суперслед не помогает}

При \(\Gamma=\operatorname{diag}(+1,-1,+1)\) ненулевые спектры на
\(\mathcal H^+\) и \(\mathcal H^-\) попарно совпадают, так что
\begin{equation}
 \Str f(D^2)=f(0)\operatorname{ind}M,
 \qquad \Str D^{2k}=0\quad(k>0).
 \label{eq:v5-supertrace-index}
\end{equation}
Для теплового ядра остаётся индекс, в согласии с формулой Маккина--Зингера
(\path{arXiv:2307.11061}), но не непрерывно меняющийся потенциал отображения
момента.

\section{Область применимости}

Теорема охватывает \(\Tr f(D^2)\), функции его спектральных моментов,
определители от сингулярных чисел и суперследы без дополнительного
ориентационного проектора. Она не охватывает условные ожидания на вершину,
относительную кривизну, вспомогательное \(D\)-поле, скрученные коммутаторы,
производные и BV--BFV-структуры, относительные модулярные и нелокальные
граничные функционалы: все они содержат дополнительные данные ориентации.

\begin{nogocriterion}[Запрет переименования]
После введения дополнительных ориентационных данных конструкция может быть
состоятельной, но её нельзя выдавать за следствие одного обычного
спектрального следа.
\end{nogocriterion}

\section{Толкование и вывод}

Точный смысл результата таков: спектр обычного нечётного оператора Дирака
помнит суммарную среднюю матрицу Грама, но забывает её разложение на входящую
и исходящую части. Отображения момента колчанов являются стандартными
симплектическими данными (\path{arXiv:0807.4734}), а спектральное действие
введено в \path{arXiv:hep-th/9606001}.

\begin{protocol}
\begin{enumerate}
 \item теорема о слепоте общего \(f(D^2)\): \textbf{доказана};
 \item следствие о знаке квартичного члена: \textbf{доказано};
 \item исправление посредством суперследа: \textbf{невозможно};
 \item обычное односледовое происхождение отображения момента:
 \textbf{направление закрыто};
 \item необычные геометрические направления: \textbf{не решены};
 \item физическое замыкание: \textbf{не достигнуто}.
\end{enumerate}
\end{protocol}

Следующая проверка --- \path{version5_nonordinary_architecture_fork_gate}.
Машинный сертификат сохранён в
\path{s2t_v5_ordinary_spectral_moment_map_no_go_results.json}.